\documentclass[12pt, titlepage]{article}
\usepackage[colorlinks=true, linkcolor=blue, urlcolor=magenta]{hyperref}

\title{Feedback Controls Lab Weekly Report \\ \normalsize{Week 9}}
\author{Aydan Vissing, Brady Schick, Gabriel Southwell, Simon Ross}
\date{\today}

\begin{document}
\maketitle

\section*{Aydan Vissing (7 hours spent)}

I took some time over the weekend to try and figure out how to use the QT editor, I think I am doing something incorrectly when it comes to it because I have not been able to figure anything out with my own experimentation. I attempted to find some videos for it but I wasn't successful at finding anything that helped. I then moved on to working on the next mini-project, battery and power storage. I spent some time researching and reading for that project. 

\section*{Brady Schick (6 hours spent)}

I spent most of the week on starting to implement the controller logic to the MATLAB simulation file from Dr. Fredette. I plan to finish that up this upcoming week if possible to get some simulation results before the presentation. The rest of the time was spent on the machine learning MP.

\section*{Gabriel Southwell (4 hours spent)}

I worked on the time of flight sensors. There are quite a few library options for using them, but none of them are very easy to use or understand. I did not make loads of progress on making them work. The one library that was kind of easy to use makes the assumption that you are only using one sensor. I tried to figure out a way to extend the library but in the end it was going to be nearly as much work as writing my own library and driver.

This coming week, I am going to keep working on it. I hope to find a library that is an easy solution, but more likely I am going to buckle in and figure out what I can do with the more powerful libraries.

\section*{Simon Ross (8 hours spent)}

This week has been less productive than I expected, but not without result. With regard to our schedule, good progress has been made on the distance sensors and we are almost ready to begin combining components for a prototype. The one remaining obstacle for this is deciding where to put each component, which will be incorporated with the PCB design.

I was planning to complete the entire PCB design over break, but I came down with the flu and was far less productive than I was hoping. The good news is that I was able to create a parts list for most of the major components that the final PCB will need (and I hope to confirm this selection with Dr. Kohl soon). As soon as I can, I will meet with Gabe to work on the final pinout of the microcontroller and any other components that I may have missed.

Beyond this, the schedule has us working on the prototype drone for the next few weeks. This means that we need to keep pressing towards the integration of our selected components, but also that we have a little time to get this figured out.

\end{document}